% =====================================================================
\section{The illness: what has to be explained}
Schizophrenia is a chronic psychiatric disorder. The latest revision of the international classification (ICD-11, in force since 2022, with countries migrating to it gradually) describes it not as a set of ready-made forms but as a profile along several independent symptom axes. Two things matter for a model: the profile differs between patients, and the course has a characteristic shape --- a long hidden preparation, a comparatively abrupt transition, and episodes thereafter.
\facts
\begin{itemize}
\item An ICD-11 diagnosis requires at least two of seven symptom groups, present most of the time for at least a month. At least one must come from the four core groups: delusions (persistent false beliefs); hallucinations, most often voices; thought disorder; and the experience of being influenced or controlled from outside. The remaining three are negative symptoms (less emotion, less drive, less speech), disorganised behaviour, and motor disturbance.\tP
\item The former forms (paranoid, hebephrenic, catatonic, simple) are abolished in ICD-11. Instead, six axes are rated separately for each patient: positive, negative, depressive, manic, motor and cognitive symptoms (memory, attention), together with the type of course: first episode, recurrent episodes, continuous.\tP
\item Catatonia --- a state of severe motor disturbance --- is now a separate diagnosis.\tP
\item Alongside schizophrenia, the same group of ``primary psychotic disorders'' contains separate diagnoses in which part of the picture is absent or short-lived: schizoaffective, schizotypal, delusional and acute transient psychotic disorder.\tP
\item \emph{Psychosis} is a state with pronounced positive symptoms; it arrives in episodes.\tU
\item Lifetime prevalence is 4.0 per 1000 (median across studies); lifetime morbid risk is 7.2 per 1000, roughly one person in a hundred and forty.\tP
\item Onset most often falls in late adolescence and early adulthood.\tU
\item Heritability is about 79\,\%.\tP\ This is the share of the variance of risk in the population attributable to genes, not the probability of passing the illness to a child. There is no single ``schizophrenia gene'': hundreds of variants are known, each contributing little.\tU
\item The first episode is preceded by a \emph{prodrome} --- months to years of non-specific change. The distribution of its length has a heavy tail: in one study the mean was 17.6 months and the median 4.\tS
\item In 43.7\,\% of patients the first psychosis developed acutely, over days.\tS
\item Drugs (antipsychotics) attenuate the positive symptoms and act little on the negative and cognitive ones.\tU
\end{itemize}
\begin{note}
The figures on prevalence, heritability and course were obtained under earlier diagnostic criteria (ICD-9, ICD-10 and the American manuals); the boundaries of the diagnosis shifted with the move to ICD-11. A long preparation followed by an abrupt transition suggests a system with a threshold, but a gradual accumulation of change with no threshold at all would produce a similar picture if symptoms appear only above some level of damage. The shape of the course is an argument, not a proof.
\end{note}
\begin{ineq}
A diagnosis is a point in a space of six axes, not the value of a single variable. A model of one mechanism may claim to explain individual axes, but not the illness as a whole. The hypothesis of this brochure bears chiefly on the cognitive axis (through PV cells and cortical rhythms) and on the positive one (through dopamine); about the depressive and manic axes it says nothing.
\end{ineq}
